% Options for packages loaded elsewhere
\PassOptionsToPackage{unicode}{hyperref}
\PassOptionsToPackage{hyphens}{url}
\documentclass[
  12pt,
]{article}
\usepackage{xcolor}
\usepackage{amsmath,amssymb}
\setcounter{secnumdepth}{-\maxdimen} % remove section numbering
\usepackage{iftex}
\ifPDFTeX
  \usepackage[T1]{fontenc}
  \usepackage[utf8]{inputenc}
  \usepackage{textcomp} % provide euro and other symbols
\else % if luatex or xetex
  \usepackage{unicode-math} % this also loads fontspec
  \defaultfontfeatures{Scale=MatchLowercase}
  \defaultfontfeatures[\rmfamily]{Ligatures=TeX,Scale=1}
\fi
\usepackage{lmodern}
\ifPDFTeX\else
  % xetex/luatex font selection
\fi
% Use upquote if available, for straight quotes in verbatim environments
\IfFileExists{upquote.sty}{\usepackage{upquote}}{}
\IfFileExists{microtype.sty}{% use microtype if available
  \usepackage[]{microtype}
  \UseMicrotypeSet[protrusion]{basicmath} % disable protrusion for tt fonts
}{}
\makeatletter
\@ifundefined{KOMAClassName}{% if non-KOMA class
  \IfFileExists{parskip.sty}{%
    \usepackage{parskip}
  }{% else
    \setlength{\parindent}{0pt}
    \setlength{\parskip}{6pt plus 2pt minus 1pt}}
}{% if KOMA class
  \KOMAoptions{parskip=half}}
\makeatother
\usepackage{longtable,booktabs,array}
\usepackage{calc} % for calculating minipage widths
% Correct order of tables after \paragraph or \subparagraph
\usepackage{etoolbox}
\makeatletter
\patchcmd\longtable{\par}{\if@noskipsec\mbox{}\fi\par}{}{}
\makeatother
% Allow footnotes in longtable head/foot
\IfFileExists{footnotehyper.sty}{\usepackage{footnotehyper}}{\usepackage{footnote}}
\makesavenoteenv{longtable}
\setlength{\emergencystretch}{3em} % prevent overfull lines
\providecommand{\tightlist}{%
  \setlength{\itemsep}{0pt}\setlength{\parskip}{0pt}}
\usepackage{bookmark}
\IfFileExists{xurl.sty}{\usepackage{xurl}}{} % add URL line breaks if available
\urlstyle{same}
\hypersetup{
  hidelinks,
  pdfcreator={LaTeX via pandoc}}

\author{SCX}
\date{2026}

\begin{document}

\section{Hostile Review Report: Line-by-Line Attack on Situs Theory}\label{hostile-review}

\begin{center}\rule{0.5\linewidth}{0.5pt}\end{center}

\textbf{Reviewer Identity}: Annals of Statistics' most demanding reviewer\\
\textbf{Review Date}: 2026-06-29\\
\textbf{Review Target}: \texttt{paper/situs\_theory/main.tex} +
\texttt{theory/self\_evolution/ppe\_rigorous\_derivation.md}\\
\textbf{Review Standard}: Impolite. Uncompromising. Not a single loophole goes unchallenged.

\begin{center}\rule{0.5\linewidth}{0.5pt}\end{center}

\subsection{0. Overall Impression}\label{overall-impression}

This paper claims to establish a ``rigorous mathematical foundation'' for Situs (physical position encoding) on top of a framework called SCX. After reading all three documents, the conclusion is: \textbf{This paper contains fundamental errors or irreparable logical leaps in the proofs of multiple core theorems. The proof of Theorem 2.4.1 is wrong---not incomplete, wrong. Theorem 1.3.1's Lipschitz constant is off by a factor of ~3.5. Theorem 2.2.1's Step 4 is a chasm that cannot be bridged with current tools. Theorem 3' ``constructive example'' actually falsifies the proposition it claims to prove. Below, itemized breakdown.}

\begin{center}\rule{0.5\linewidth}{0.5pt}\end{center}

\subsection{1. Notation Consistency: main paper vs ppe\_rigorous\_derivation.md}\label{notation-consistency}

\subsubsection{1.1 Inconsistent Theorem Numbering}\label{inconsistent-theorem-numbering}

\begin{longtable}[]{@{}p{0.10\linewidth}p{0.22\linewidth}p{0.55\linewidth}p{0.10\linewidth}@{}}
\toprule\noalign{}
\begin{minipage}[b]{\linewidth}\raggedright
Location
\end{minipage} & \begin{minipage}[b]{\linewidth}\raggedright
main.tex Number
\end{minipage} & \begin{minipage}[b]{\linewidth}\raggedright
ppe\_rigorous\_derivation.md Number
\end{minipage} & \begin{minipage}[b]{\linewidth}\raggedright
Diff
\end{minipage} \\
\midrule\noalign{}
\endhead
\bottomrule\noalign{}
\endlastfoot
Learnable PE Destruction Theorem & Theorem 4.2 (line 1048) & Theorem 4.1 (\S4.2.3, line 851) &
\textbf{Number off by 1} \\
Encoding Imperfection Upper Bound & Proposition 3.1 (line 864) & Proposition 3.1 (\S3.1.3, line 690) &
Consistent \\
Theorem 2' & Theorem 2' (line 881) & Theorem 2' (\S3.2, line 702) &
Consistent \\
\end{longtable}

\textbf{Attack}: The Theorem 4.1 vs Theorem 4.2 numbering gap implies a ``vanished Theorem 4.1'' exists in the paper body. If Theorem 4.1 is Proposition 4.1 (fixed PE preserves Theorem 3), then why isn't Proposition 4.1 called Theorem 4.1? This is not a typo---it is a break in the logical chain: Proposition 4.1 uses a trivial deterministic function argument that doesn't qualify as a theorem; yet Theorem 4.2's ``destructive'' result is the true theorem-level contribution. The numbering confusion exposes the authors' own confusion about contribution importance.

\subsubsection{1.2 Split Personality of the Symbol Naming System}\label{split-personality-symbols}

The paper body uses \texttt{\textbackslash{}Situs}, \texttt{\textbackslash{}Sstates}, \texttt{\textbackslash{}Ppos}, \texttt{\textbackslash{}PE}, \texttt{d\_\{\textbackslash{}pe\}}, while ppe\_rigorous\_derivation.md uses \texttt{PPE}, \texttt{\textbackslash{}mathcal\{S\}}, \texttt{\textbackslash{}mathcal\{P\}}, \texttt{\textbackslash{}text\{PE\}}, \texttt{d\_\{\textbackslash{}text\{pe\}\}}. The paper claims these are ``consistent'' with ppe\_rigorous\_derivation.md (line 1343), yet they differ in every symbol. This is not a LaTeX vs Markdown formatting difference---the paper uses \texttt{\textbackslash{}Ppos} to denote position space, while md uses \texttt{\textbackslash{}mathcal\{P\}}. The authors demonstrate inconsistency in the very same sentence that claims consistency.

\subsubsection{1.3 Two Versions of Theorem 2.3.1}\label{two-versions-theorem-231}

\textbf{Paper version} (line 697-701):
\[\delta_s^{\PE} \leq \min\!\left(1 - p_{\text{clean}, s},\, \sqrt{\frac{1}{2} D_{\KL}(P_{\PE|S,Y} \| P_{\PE|S})}\right) + \delta_s^{\text{variance}}\]

\textbf{md version} (line 512-513):
\[\delta_s^{\text{PE}} \leq \min\left(1, \sqrt{\frac{1}{2} \cdot I(P; \text{PE}(P) \mid S, Y)}\right) + \text{Bernstein correction term}\]

These are \textbf{not} the same upper bound. The first uses \(1-p_{\text{clean},s}\) and KL divergence (pointwise conditional), the second uses \(1\) and mutual information (expectation form). The md version uses ``Bernstein correction'' (implying \(O(1/M)\) convergence), while the paper uses \(\delta_s^{\text{variance}} = O(1/\sqrt{M})\) (implying CLT convergence). Their scaling laws under finite \(M\) are completely different---one is \(1/M\), the other is \(1/\sqrt{M}\).

\begin{center}\rule{0.5\linewidth}{0.5pt}\end{center}

\subsection{2. Assumption Verification: Do A1--A6 Still Hold Under Situs?}\label{assumption-verification}

The CC audit report implicitly contains 6 core assumptions for the original SCX theorems. Let us check each after adding Situs:

\subsubsection{A1: Expert Independence (\(v_m \indep v_{m'} \mid s\))}\label{a1-expert-independence}

\textbf{Original state}: \(M\) experts' votes on the same state atom \(s\) are i.i.d. Bernoulli.

\textbf{After adding Situs}: All experts receive the same augmented input \(h_i = \phi(s_i) + \PE(p_i)\). The PE term is a shared deterministic signal---if PE carries information about the label, then all experts' inputs contain the same ``hint.'' This \textbf{weakens} the conditional independence assumption: \(\Cov(v_m, v_{m'} \mid s, p) \neq 0\). Strictly speaking, the Chernoff-Hoeffding bound requires conditional independence given the state, and now PE injects a common signal shared by all experts. \textbf{This breaks the strictness of the i.i.d. assumption.}

Attack severity: \textbf{Medium}. When \(M\) is large, the correlation introduced by this shared signal may loosen the Chernoff bound by a constant factor (similar to \(\beta\)-mixing correction), but the paper does not mention this issue---it pretends the Chernoff-Hoeffding structure is ``unchanged.''

\subsubsection{A2: Positive Detection Margin (\(\Delta_s > 0\))}\label{a2-positive-detection-margin}

\textbf{Original state}: \(\Delta_s = p_{\text{noisy},s} - p_{\text{clean},s} > 0\).

\textbf{After adding Situs}: We need \(\Delta_s^{\Situs} = \Delta_s + \delta_s^{\PE} > 0\). But \(\delta_s^{\PE} \in [-\Delta_s, 1-p_{\text{clean},s}]\). \textbf{Situs can reduce the detection margin!} If \(\delta_s^{\PE} < 0\) and \(|\delta_s^{\PE}| > \Delta_s\), then \(\Delta_s^{\Situs} < 0\)---Situs reverses the detection direction. Theorem 2.2.1 only gives a sufficient condition for the existence of some \(s\) and some PE such that \(\delta_s^{\PE} > 0\), but \textbf{does not rule out the possibility that \(\delta_s^{\PE} < 0\) for other \(s\)}.

Attack severity: \textbf{High}. The paper's corrected Theorem 1 replaces \(\Delta_s\) with \(\Delta_s + \delta_s^{\PE}\) for all states, but \(\delta_s^{\PE}\) can be negative. For some \(s\), \(\exp(-2M(\Delta_s + \delta_s^{\PE})^2)\) may be larger than \(\exp(-2M\Delta_s^2)\), i.e., the \(F_1\) lower bound \textbf{gets worse}. This is not the ``systematic enhancement'' the paper claims.

\subsubsection{A3: State Space Discreteness}\label{a3-state-space-discreteness}

\textbf{Original state}: \(\mathcal{S} = \{s_1, \ldots, s_N\}\) is a finite discrete set.

\textbf{After adding Situs}: The encoding space \(\R^d\) is continuous, and \(h_i = \phi(s_i) + \PE(p_i)\) varies continuously in \(\R^d\). Even if \(s_i\) is discrete, \(h_i\) is no longer. This means ``equivalence classes'' among state atoms are no longer strict---two atoms originally mapped to the same state \(s\) will receive different \(h_i\) if their physical positions differ. This makes the definition of probability \(\rho_s\) for ``aggregation by state \(s\)'' ambiguous---originally aggregated by state atom type, now must be aggregated jointly by \((s, p)\).

Attack severity: \textbf{Low-Medium}. Technically fixable (define distribution on the joint space), but the paper does not discuss this subtlety.

\subsubsection{A4: Bounded Random Variables (Chernoff-Hoeffding Applicability)}\label{a4-bounded-random-variables}

\textbf{Original state}: \(v_m \in \{0,1\}\) is bounded.

\textbf{After adding Situs}: Votes remain in \(\{0,1\}\), Chernoff-Hoeffding applicability is unchanged. This assumption \textbf{formally holds}. But as noted in A1, the independence condition is weakened.

Attack severity: \textbf{Low} (formally preserved, substantively implicated by A1).

\subsubsection{A5--A6: Bayes Error Rate and Fano Condition}\label{a5-a6-bayes-fano}

These involve internal constructions of the original Theorem 2 and Theorem 3, not explicitly numbered in the CC audit report. Under Situs:
- Bayes error rate in the augmented feature space requires extra care (augmented space introduces continuous dimensions)
- The \(\log|\Y|\) term in Fano's inequality is unchanged, but conditional entropy computation involves continuous-discrete mixtures

Attack severity: \textbf{Low}. These are standard information-theoretic technical details that can be handled, but the paper does not handle them.

\subsubsection{Assumption Verification Conclusion}\label{assumption-verification-conclusion}

\textbf{Of the 6 assumptions claimed by the paper, at least A1 (independence) and A2 (positive margin) are significantly affected under Situs. The paper does not verify these assumptions one by one yet claims the Chernoff structure is ``unchanged''---this is dishonest.}

\begin{center}\rule{0.5\linewidth}{0.5pt}\end{center}

\subsection{3. Proof Completeness: Logical Leaps, Undeclared Assumptions, Erroneous Reasoning}\label{proof-completeness}

\subsubsection{3.1 Theorem 2.2.1 Step 4: Bayes Optimal $\to$ Actual Expert Generalization (Fatal)}\label{theorem-221-step4}

The paper body line 659-661:
\textgreater{} Step 4 (from Bayes optimal to actual experts): Assuming actual experts consistently approximate Bayes optimal (extracting information from augmented features in the same direction as Bayes optimal), the above inequality sign holds in expectation.

The paper itself annotates this as \texttt{[heuristic]}. But this is not heuristic---\textbf{this is an unbridgeable chasm}. The issues are:

\textbf{(a)} Between the behavior of the Bayes optimal classifier given \((s, \PE(p))\) and the behavior of any actual trained neural network expert \(E_m\), there is no finite-sample convergence guarantee.

\textbf{(b)} What does ``consistently approximate Bayes optimal'' mean? If the experts are neural networks trained by SGD, their convergence points are some stationary points in the loss landscape, not necessarily Bayes optimal. Even if they were, it is in the \textbf{population distribution} sense, not guaranteed to have the correct sign of \(\delta_s^{\PE}\) for every specific state \(s\).

\textbf{(c)} The paper claims \(p_{\text{clean},s}^{\Situs} \leq p_{\text{clean},s}\) (disagreement on clean samples does not increase). A Bayes optimal classifier, given more information, indeed cannot get worse (data processing inequality guarantees this). But actual experts may perform \textbf{worse} on clean samples because PE introduces \textbf{spurious correlations}---especially when training data is limited, PE may learn a pattern that works on the training set but not on the test set. In this case \(p_{\text{clean},s}^{\Situs} > p_{\text{clean},s}\), and \(\delta_s^{\PE}\) may be negative.

\textbf{(d)} The entire conclusion of Theorem 2.2.1---``there exists PE and dimension such that \(\delta_s^{\PE} > 0\)''---without Step 4, holds only for the Bayes optimal classifier. For the neural network experts in the actual SCX framework, this is an \textbf{unproven conjecture}.

\textbf{Verdict}: Theorem 2.2.1 is the paper's most important information-theoretic result (connecting \(I(Y;P|S) > 0\) and detection improvement). Without Step 4, it only proves that Bayes optimal classifiers can benefit from positional information---which is trivial. Extending this result to actual experts in the SCX framework is a \textbf{logical leap}.

\subsubsection{3.2 Theorem 2.4.1 Step 2: Fano's Inequality Direction (Fatal)}\label{theorem-241-fano}

This is the most serious error I discovered. The paper claims Fano gives a \textbf{lower bound} for \(\delta_s^{\PE}\). But examining the derivation carefully:

\textbf{Step 1}: Fano's inequality in standard form (correct):
\[H(Y \mid S, \PE(P)) \leq H_2(P_e) + P_e \cdot \log(|\Y| - 1)\]

From this (paper line 750-752):
\[P_e \geq \frac{H(Y \mid S, \PE(P)) - 1}{\log |\Y|}\]

This is a correct \textbf{lower bound}.

\textbf{Step 3} (paper line 763-766, md line 601-607):
\[P_e - P_e^{\Situs} = \frac{H(Y|S) - H(Y|S, \PE(P))}{\log |\Y|} \quad (\text{taking Fano's equality})\]

\textbf{This is a fundamental error.} The logical chain is:
1. \(P_e \geq \frac{H(Y|S)-1}{\log|\Y|}\) (Fano)
2. \(P_e^{\Situs} \geq \frac{H(Y|S,\PE(P))-1}{\log|\Y|}\) (Fano)
3. If both inequalities achieve equality $\to$ \(P_e - P_e^{\Situs} = \frac{H(Y|S) - H(Y|S,\PE(P))}{\log|\Y|}\)

But from the \textbf{inequalities} in (1) and (2) to the \textbf{equality} in (3), we need (i) Fano's inequality to achieve equality in both worlds, and (ii) ``the difference of two Fano equalities = the difference of error rates.''

The problem is:
- Fano's inequality achieves equality only for specific distributions (noise uniformly distributed across wrong classes, and \(H_2(P_e)\) is binary entropy), which does not hold in general.
- More importantly, even if both Fano bounds achieve equality, \(a \geq A\) and \(b \geq B\) \textbf{cannot imply} \(a-b = A-B\) or \(a-b \geq A-B\).

\textbf{The essence of the directional error}: The paper attempts to derive a \textbf{lower bound} on error rate improvement from Fano's \textbf{upper bound} (conditional entropy). Fano's inequality has the form \(H(Y|X) \leq f(P_e)\), giving a lower bound on \(P_e\). But \(P_e - P_e^{\Situs}\) as a difference cannot be directly derived from two lower bounds on \(P_e\)---the direction of the difference is not controlled by the individual lower bounds.

\textbf{Verdict}: \textbf{The proof of Theorem 2.4.1 is wrong.} Not incomplete, not heuristic---it is a logical fallacy of deriving an equality from inequalities. The word ``lower bound'' in the theorem statement misleads the reader into thinking this is a rigorous information-theoretic bound. In reality it is at best a heuristic guess.

\subsubsection{3.3 Theorem 2.2.1 Step 3's ``Clean Sample Disagreement Does Not Increase'' Argument}\label{clean-sample-disagreement}

Paper line 655: \textgreater{} For clean samples, the additional position information further confirms the correct label $\to$ disagreement probability does not increase: \(p_{\text{clean}, s}^{\Situs} \leq p_{\text{clean}, s}\).

This relies on an implicit assumption: PE's information is ``consistent'' with the correct label rather than ``conflicting.'' But:
- PE is derived from physical position, which may have \textbf{spurious correlations} with the label (correlated in the training distribution but not in the target distribution).
- PE may amplify some systematic bias of the experts---if all experts share a common blind spot, additional PE information may make them collectively more confident in being wrong on that blind spot.
- The paper itself acknowledges this risk: in \S6.3's ``one-vote veto condition,'' F5 points out that under ``random or pseudo-random configurations + finite samples,'' \(\delta_s^{\PE} > 0\) on the training set may become \(\delta_s^{\PE} < 0\) on the test set.

\textbf{Attack}: Step 3's argument assumes PE-provided information is always ``beneficial''---it increases disagreement on noisy samples but not on clean samples. This is equivalent to \textbf{assuming the conclusion}: PE somehow distinguishes between noise and clean samples. If PE's information is a pure additional feature (correlated but not causal), it may increase disagreement on both types of samples by the same magnitude, yielding \(\delta_s^{\PE} = 0\).

\subsubsection{3.4 Theorem 2.5.1: Illegitimate Use of Lipschitz Continuity Assumption}\label{theorem-251-lipschitz}

Theorem 2.5.1 claims:
\[|\delta_i^{\PE} - \delta_j^{\PE}| \leq 2 \cdot L_E \cdot L_{\PE} \cdot d_{\Ppos}(p_i, p_j)\]

The proof uses (md line 640):
\[|p_{\text{clean}, i}^{\text{PPE}} - p_{\text{clean}, j}^{\text{PPE}}| \leq L_E \cdot \|h_i - h_j\|\]

But \(p_{\text{clean}, i}^{\text{PPE}} = \mathbb{E}[\mathbf{1}[E_m(h_i) \neq y]]\) is the \textbf{expectation of an indicator function}. Even if \(E_m\)'s output is \(L_E\)-Lipschitz (in some sense), the indicator function \(\mathbf{1}[\cdot \neq y]\) is \textbf{discontinuous}---it jumps at the decision boundary. The \(L_E\)-Lipschitz assumption typically applies to the network's output logits or softmax probabilities, not to the expectation of 0/1 decisions.

For this bound to hold, additional assumptions are needed:
- The experts' decision boundaries are sufficiently far from both \(h_i\) and \(h_j\) (margin condition), or
- The experts output soft probabilities and the expectation of the indicator function happens to be Lipschitz (requiring the experts' predicted probabilities to be Lipschitz in the input, and the probability of threshold crossing to be a zero-measure set)

The paper provides none of these conditions. \textbf{This is not a technical detail---it is crucial for whether the bound is non-vacuous.}

\begin{center}\rule{0.5\linewidth}{0.5pt}\end{center}

\subsection{4. Sign Errors: More Than One}\label{sign-errors}

\subsubsection{4.1 Sign Error in CC Audit Report (Fixed, Confirmed)}\label{sign-error-cc}

The CC audit report's Proposition 2.1 indeed has a sign error. Its
\(\delta_s^{\PE} = (p_{\text{clean}}^{\text{PPE}} - p_{\text{clean}}) - (p_{\text{noisy}}^{\text{PPE}} - p_{\text{noisy}})\)
has the physical interpretation: encoding is helpful only when the consensus gain on clean samples exceeds the ``confusion gain'' on noisy samples. This contradicts the intuition that ``encoding should increase the noise detection margin.'' The paper's correction (swapping the order of the two terms) is correct.

\textbf{However}: The CC audit report's internal notation is itself inconsistent. \S0.2 defines \(p_{\text{clean}}\) as disagreement probability (\(<0.5\)), but \S2.3's proof contains \(|p_{\text{clean}} + p_{\text{noisy}} - 1|\), a formula that treats \(p_{\text{clean}}\) as consensus probability. So the CC report's ``sign error'' may not be simply a reversed sign---it is a deeper \textbf{internal split in the notation system}: \S0.2 and \S2.3 use completely different semantics for \(p\).

The paper fixes this sign error but \textbf{does not point out the CC report's internal notation split}---it simply flips the sign and claims correction. In reality the CC report's problem is more severe.

\subsubsection{4.2 Theorem 1.3.1: Massive Overestimation of the Lipschitz Constant (New Derivation Error)}\label{theorem-131-lipschitz-overestimation}

This is not a sign error---it is a magnitude error. The paper claims:
\[L_{\PE}^{\text{rot}} = 2\sqrt{\alpha^2 + \beta^2 + \gamma^2}\]

The actual Lipschitz constant is \(\max(\alpha, \beta, \gamma)\). Off by a factor of \(2\sqrt{3} \approx 3.46\) (when \(\alpha=\beta=\gamma\)).

\textbf{Derivation}: For \(\mathbf{e}_0 = (1,0,0,1,0,0,\ldots)^T\) (first dimension of each rotation plane is 1),
\[\|\PE_{\text{rot}}(\mathbf{p}) - \PE_{\text{rot}}(\mathbf{q})\|^2 = 4\sin^2(\alpha\Delta x/2) + 4\sin^2(\beta\Delta y/2) + 4\sin^2(\gamma\Delta z/2)\]

As \(\|\Delta\mathbf{p}\| \to 0\):
\[\|\PE_{\text{rot}}(\mathbf{p}) - \PE_{\text{rot}}(\mathbf{q})\| \approx \sqrt{\alpha^2\Delta x^2 + \beta^2\Delta y^2 + \gamma^2\Delta z^2}\]

By Cauchy-Schwarz:
\[\sqrt{\alpha^2\Delta x^2 + \beta^2\Delta y^2 + \gamma^2\Delta z^2} \leq \max(\alpha, \beta, \gamma) \cdot \|\Delta\mathbf{p}\|_2\]

\textbf{The true Lipschitz constant is \(\max(\alpha, \beta, \gamma)\)}, and this bound is tight (take \(\Delta\mathbf{p}\) along the direction of the largest frequency parameter).

The error in the paper's proof stems from the triangle inequality decomposition in Step 2 (md line 300-317). The authors correctly obtain:
\[\|\mathbf{R}(\mathbf{p}) - \mathbf{R}(\mathbf{q})\|_F \leq 2\alpha|\Delta x| + 2\beta|\Delta y| + 2\gamma|\Delta z|\]

But the three difference matrices act on \textbf{mutually disjoint subspaces} (dimension pairs (0,1), (2,3), (4,5)). For the sum of matrices on disjoint subspaces \(D = D_x + D_y + D_z\):
\[\|D\|_F^2 = \|D_x\|_F^2 + \|D_y\|_F^2 + \|D_z\|_F^2\]

The triangle inequality \(\|D\|_F \leq \|D_x\|_F + \|D_y\|_F + \|D_z\|_F\)
can be loose by a factor of up to \(\sqrt{3}\) in the worst case (three equal-norm terms). Then the authors introduce another \(\sqrt{3}\) factor of looseness via Cauchy-Schwarz (\(\ell_1 \to \ell_2\) conversion).

\textbf{The product of these two loosenings yields the final constant being off by approximately 3.46$\times$}. Worse, the authors noticed this issue (Note 1.3.1/md note 1.3.1 acknowledges that the effective single-direction constant is only \(\alpha\), not \(2\alpha\)), yet still write the loose bound as the ``exact constant'' in the theorem statement.

\textbf{Verdict}: Theorem 1.3.1's claimed ``exact Lipschitz constant'' is neither exact nor tight. The true exact Lipschitz constant (using \(\|\cdot\|_2\) vector norm and \(\|\cdot\|_2\) position norm) is \(\max(\alpha, \beta, \gamma)\).

\subsubsection{4.3 Theorem 1.2.1: Missing Normalization Factor}\label{missing-normalization}

The encoding kernel definition (Definition 1.2.1) is:
\[\langle \PE_{\text{scalar}}(p), \PE_{\text{scalar}}(q) \rangle = \sum_{j=0}^{d/2-1} \cos(2\pi(p-q)/\lambda_j) = K_{\PE}(\Delta)\]

At \(\Delta = 0\), \(K_{\PE}(0) = d/2\). The target kernel has \(k(0) = 1\).

As \(d \to \infty\), the encoding kernel \textbf{diverges} to infinity rather than converging to the target kernel. For the encoding kernel to approximate the target kernel, we must normalize:
\[\frac{2}{d} K_{\PE}(\Delta) \to \int_0^\infty S(\omega)\cos(\omega\Delta)d\omega = k(\Delta)\]

But Theorem 1.2.1's statement and proof \textbf{completely omit} this \(2/d\) factor. The theorem's claim that ``the encoding kernel optimally approximates \(k(\Delta)\) in \(L^2([0,L])\)'' is mathematically false---\(K_{\PE}\) and \(k\) have different magnitudes (\(O(d)\) vs \(O(1)\)).

\subsubsection{4.4 Theorem 2.3.1: Legitimacy of Pinsker Application}\label{pinsker-legitimacy}

Theorem 2.3.1 uses Pinsker's inequality to connect \(\delta_s^{\PE}\) with KL divergence. But there is a subtle issue with Pinsker's inequality:

Pinsker's bound is: \(|P(A) - Q(A)| \leq \sqrt{\frac{1}{2}D_{\KL}(P\|Q)}\), where \(P\) and \(Q\) are probability measures on the \textbf{same measurable space} and \(A\) is an event in that space.

In the paper's application, \(P\) and \(Q\) are the conditional distributions of \(\PE(P)\) given \(S\) (clean vs mixed), but the quantity being bounded, \(p_{\text{clean},s}^{\Situs} - p_{\text{clean},s}\), involves \textbf{the expert's output event}, which is not in the space of \(\PE(P)\).

Strictly speaking, the correct argument would be:
\[p_{\text{clean},s}^{\Situs} = \int P(E_m(\phi(s)+z) \neq y) \cdot dP_{\PE(P)|S=s,\text{clean}}(z)\]
\[p_{\text{clean},s} = P(E_m(\phi(s)) \neq y)\]

The second quantity does not even involve the distribution of \(\PE(P)\). The paper compares \(p_{\text{clean},s}^{\Situs}\) and \(p_{\text{clean},s}\), but Pinsker only constrains the change in the distribution of \(\PE(P)\), not directly the change in output probability---unless we additionally assume the expert's output is Lipschitz in its dependence on PE (returning to the problem of Theorem 2.5.1).

\textbf{However, this may be fixable}: If we compare \(p_{\text{clean},s}^{\Situs}\) and \(p_{\text{clean},s}^{\text{avg}} = \int P(E_m(\phi(s)+z) \neq y) \cdot dP_{\PE(P)|S=s}(z)\) (i.e., clean disagreement rate averaged over PE), Pinsker can be legitimately applied. But the paper does not make this decomposition.

\begin{center}\rule{0.5\linewidth}{0.5pt}\end{center}

\subsection{5. Tightness: Claims of Tightness But None Are Tight}\label{tightness}

\subsubsection{5.1 Theorem 1.2.3: Scalar Sinusoidal Lipschitz Constant}\label{theorem-123-tight}

Claimed to be ``tight'' (line 414, md line 203). Each \((\sin,\cos)\) pair's derivative indeed satisfies \(\cos^2 + \sin^2 = 1\) identically for all \(p\), so for any two close points, the local rate of change of this pair indeed reaches \(\omega_j\). For the whole vector, all pairs act independently in orthogonal subspaces, so the Lipschitz constant is indeed \(\sqrt{\sum_j \omega_j^2 \cdot 2}\) = the claimed value. \textbf{This bound is indeed tight.} $\checkmark$

\subsubsection{5.2 Theorem 1.3.1: 3D Rotation Lipschitz Constant}\label{theorem-131-not-tight}

Claimed to be ``worst-case'' tight (line 519, md line 337). As shown in \S4.2, \textbf{not tight}---off by a factor of ~3.46. The true tight Lipschitz constant is \(\max(\alpha, \beta, \gamma)\). $\times$

\subsubsection{5.3 Theorem 1.2.2: Approximation Error Bound}\label{theorem-122-approximation}

Claims \(O(1/d)\) with constant \(\frac{2\pi^2\xi}{d} \cdot \frac{L}{\xi} = \frac{2\pi^2 L}{d}\). \textbf{No proof.} Moreover, \(\xi\) appears in both numerator and denominator and cancels, indicating the constant \(\frac{2\pi^2 L}{d}\) is independent of the physical correlation length \(\xi\)---this is extremely suspicious. How can a kernel's approximation quality be independent of the kernel's width?

Actual derivation (what I did): For Riemann-Stieltjes sum approximation of \(\cos(\omega\Delta)\) on interval \([0,L]\) using quantile sampling, the error is jointly controlled by truncated high-frequency components and the coarsest low-frequency resolution. With fixed \(d\), the error should vary with \(\xi\)---small \(\xi\) (short-range correlation) means more high-frequency components, harder to approximate. The paper's claimed constant being independent of \(\xi\) is physically absurd.

\subsubsection{5.4 Theorem 2' Bound}\label{theorem-2-prime-bound}

Claims \(F_{1,\text{SCX}}^{\Situs} \leq F_{1,\text{base}} + C_F \cdot \sqrt{(\delta + 2\varepsilon_{\PE}/C_F^2)/2}\).

Note that when \(\varepsilon_{\PE} \to I(Y;P|X)\) (complete encoding failure):
\[\delta_{\text{PPE}} = \delta + \frac{2I(Y;P|X)}{C_F^2}\]

The upper bound is looser than the original Theorem 2 by a factor of \(\sqrt{\delta + 2I/C_F^2} / \sqrt{\delta}\). When \(I(Y;P|X)\) is large (position has lots of information but encoding completely loses it), the upper bound can become arbitrarily loose. This is certainly not a ``tight'' bound---it is a \textbf{non-vacuous but potentially very loose} upper bound. The paper does not claim it is tight (this is honest), but the abstract says ``corrected the weak feature failure upper bound'' suggesting some precision, when in fact it merely added a positive term to an already loose bound.

\begin{center}\rule{0.5\linewidth}{0.5pt}\end{center}

\subsection{6. Theorem 3' Constructive Example: Shooting Itself in the Foot}\label{theorem-3-prime-example}

\subsubsection{6.1 Claim vs Reality Contradiction}\label{claim-vs-reality}

\S4.3 (paper line 1096-1137) is titled ``Constructive Minimal Example.'' What does this ``constructive minimal example'' demonstrate?

\begin{enumerate}
\def\labelenumi{\arabic{enumi}.}
\tightlist
\item
  Constructs two worlds \(W_A, W_B\), verifies \(P_{W_A}(X,Y) = P_{W_B}(X,Y)\) $\checkmark$
\item
  Introduces learnable PE \(\PE_\theta(p) = \theta \cdot p\)
\item
  Computes \(\mathbb{E}_{W_A}[Y|X] = \mathbb{E}_{W_B}[Y|X] = 0.6 \cdot \text{sign}(x)\)---\textbf{numerically identical}
\item
  Conclusion: \textbf{Even learnable PE cannot distinguish the two worlds} (paper line 1131-1132 original text: ``Therefore even learnable PE cannot distinguish the two worlds. This precisely confirms the essence of Theorem 3'')
\end{enumerate}

\textbf{This is a ``constructive counterexample''---it constructs an example where learnable PE cannot distinguish, precisely proving Theorem 3's robustness rather than its fragility!}

\subsubsection{6.2 Hand-Waving ``Conditions That Truly Distinguish''}\label{hand-waving}

After proving that learnable PE cannot distinguish, the paper (line 1134-1136) waves its hand and says:
\textgreater{} ``However, when the auxiliary loss involves an unobserved physical quantity \(Z\)---for example in \(W_A\) noise labels cause \(|Z_{\text{pred}} - Z_{\text{DFT}}|\) to be large (physically inconsistent), while in \(W_B\) this difference is small (difficult but physically consistent)---then the optimal parameters of the learnable PE differ between the two worlds, and unidentifiability is broken.''

This is \textbf{not a constructive proof}---it is a verbal description of an imagined scenario. No specific definition of \(Z\) is given, no specific form of \(\mathcal{L}_{\text{aux}}\) is given, no verification that the two worlds indeed yield different optimal points. The entire \S4.3 promised a ``constructive minimal example'' but only delivered an \textbf{empty claim}.

\subsubsection{6.3 Final Verdict on Theorem 3'}\label{final-verdict-theorem-3}

Theorem 3' itself is a reasonable conditional theorem statement (given premises (A)(B), unidentifiability is preserved), but:
- The ``distinguishability'' conclusion when premise (A) fails \textbf{has no constructive proof to support it}
- The ``constructive minimal example'' is actually a counterexample to the counterexample---i.e., it demonstrates the theorem's robustness
- The actual mechanism for distinguishing noisy from hard samples (multi-task learning + physical consistency) is not formalized, not proven, not constructed

\textbf{Theorem 3''s ``destruction'' part is an unproven conjecture.}

\begin{center}\rule{0.5\linewidth}{0.5pt}\end{center}

\subsection{7. Fatal Defect Summary}\label{fatal-defect-summary}

\subsubsection{Fatal Defect \#1: Theorem 2.4.1's Fano Derivation}\label{fatal-1}

\textbf{Nature}: Logical error (misuse of inequality direction)\\
\textbf{Location}: main.tex \S3.4 (line 746-776), md \S2.4 (line 557-617)\\
\textbf{Specific error}: Deriving a lower bound on error rate difference from two Fano lower bounds. \(a \geq A\) and \(b \geq B\) cannot imply \(a-b = A-B\) or \(a-b \geq A-B\).\\
\textbf{Fixability}: Not fixable. Requires a completely different proof strategy. Possible directions include strong Fano or inverse Fano, but these require specific assumptions about the distribution.\\
\textbf{Consequence}: Theorem 2.4.1 (Fano lower bound) cannot be called a theorem---it is at best a heuristic estimate.

\subsubsection{Fatal Defect \#2: Theorem 1.3.1's Lipschitz Constant}\label{fatal-2}

\textbf{Nature}: Constant factor overestimated by ~3.46$\times$\\
\textbf{Location}: main.tex \S2.3 (line 483-492), md \S1.3.4 (line 269-276)\\
\textbf{Specific error}: Double looseness from triangle inequality + Cauchy-Schwarz when subspace orthogonality should simplify to direct Frobenius norm computation.\\
\textbf{Fixability}: Fixable. The correct Lipschitz constant is \(\max(\alpha, \beta, \gamma)\).\\
\textbf{Consequence}: Theorem 1.4.1 (Unified Lipschitz) and Theorem 2.5.1 (Spatial Smoothness) both use this overestimated constant. The spatial smoothness bound is actually ~3.5$\times$ tighter than claimed---this is good news, but it shows the paper's analysis is not precise.

\subsubsection{Fatal Defect \#3: Theorem 2.2.1's Bayes $\to$ Actual Expert Chasm}\label{fatal-3}

\textbf{Nature}: Unbridgeable logical leap\\
\textbf{Location}: main.tex line 659-661, md line 484\\
\textbf{Specific error}: Bayes optimal classifier behavior cannot be generalized to actual trained neural network experts. No finite-sample guarantee.\\
\textbf{Fixability}: Possibly patchable by assuming the expert model class has some uniform approximation property, but this would fundamentally change the theorem's nature (from information theory to learning theory).\\
\textbf{Consequence}: Theorem 2.2.1 (the paper's most important connection theorem) degenerates after removing Step 4 to: \(I(Y;P|S) > 0\) implies Bayes optimal classifiers perform better when using \(P\) information. This is trivial, not a contribution.

\subsubsection{Fatal Defect \#4: Theorem 1.2.1 Missing Normalization Factor}\label{fatal-4}

\textbf{Nature}: Derivation omission\\
\textbf{Location}: main.tex \S2.2.2, md \S1.2.2\\
\textbf{Specific error}: \(K_{\PE}(0) = d/2\) while \(k(0) = 1\). Without normalization, convergence is impossible.\\
\textbf{Fixability}: Fixable. Add \(1/\sqrt{d/2}\) factor to the encoding definition or \(2/d\) normalization to the inner product. But all subsequent constants (Lipschitz, error bounds) need corresponding adjustment.\\
\textbf{Consequence}: Theorem 1.2.1's ``optimal approximation'' statement and Theorem 1.2.2's approximation error bound are wrong without normalization.

\subsubsection{Semi-Fatal Defect \#5: Theorem 3''s ``Destruction'' Direction Has No Constructive Proof}\label{semi-fatal-5}

\textbf{Nature}: Unsupported claim\\
\textbf{Location}: main.tex \S4.3, md \S4.3\\
\textbf{Specific error}: The ``constructive minimal example'' actually proves Theorem 3's robustness.\\
\textbf{Fixability}: Requires a genuine example showing that learnable PE can produce different parameters on observationally equivalent data. This requires an auxiliary loss involving unobserved variables---such an example is possible but \textbf{has not been constructed}.\\
\textbf{Consequence}: Theorem 3''s ``if premise (A) does not hold'' part is currently an unproven conjecture.

\begin{center}\rule{0.5\linewidth}{0.5pt}\end{center}

\subsection{8. Additional Issues: Small But Real}\label{additional-issues}

\subsubsection{8.1 Theorem 2.3.1: Mysterious Origin of \(\min(1-p_{\text{clean},s}, \ldots)\)}\label{mysterious-min}

Why does the upper bound contain a \(1-p_{\text{clean},s}\) term? The derivation (line 708-719) starts from Pinsker + triangle inequality to get:
\[\delta_s^{\PE} \leq \sqrt{\frac{1}{2}D_{\KL}^{(1)}} + \sqrt{\frac{1}{2}D_{\KL}^{(2)}}\]

Then suddenly becomes \(\min(1-p_{\text{clean},s}, \ldots)\). Where did \(1-p_{\text{clean},s}\) come from? It looks like \(\delta_s^{\PE}\)'s range upper bound (from \(\delta_s^{\PE} \in [-\Delta_s, 1-p_{\text{clean},s}]\)), but what is the mathematical justification for taking \(\min\) of the range upper bound and the Pinsker upper bound? If Pinsker gives 0.3 and the range upper bound gives 0.7, why is the bound 0.3 rather than 0.7? This requires arguing that the Pinsker bound \(\leq\) the range upper bound always holds, or explaining why both constitute valid upper bounds. The introduction of this \(\min\) currently has no logic.

\subsubsection{8.2 Theorem 2.4.1: Unit Ambiguity of \(\log 2\)}\label{log-2-ambiguity}

The paper's Fano lower bound contains \(\frac{2\rho I - \log 2}{\log|\Y|}\). \(I(\cdot;\cdot)\) is typically in nats (if using natural logarithm), so what unit is \(\log 2\) in? If uniformly using nats, \(H_2(P_e) \leq \ln 2\) nats; if uniformly using bits, \(H_2(P_e) \leq 1\) bit, no \(\log 2\). The appearance of \(\log 2\) indicates the paper switches between nats and bits without declaration. In an information theory paper, this is unacceptably sloppy.

\subsubsection{8.3 Theorem 1.4.1 Is Pure Packaging}\label{theorem-141-packaging}

Theorem 1.4.1 collects two Lipschitz constants in a table and claims it is a ``Unified Lipschitz Continuity Theorem.'' It contains no new mathematical content---it simply defines the union of two constants. This is not a theorem, it is a definition or summary table.

\subsubsection{8.4 Proposition 3.1's (Dimension-Information Tradeoff) Exponential Decay Claim}\label{exponential-decay-claim}

\[\varepsilon_{\PE} \leq I(Y; P \mid X) \cdot \exp(-d/d_0)\]

Annotated as \texttt{[heuristic]}. Exponential decay for Fourier series approximation of analytic kernels (Cauchy) holds under certain smoothness conditions, but what is being approximated here is the kernel \(k(\Delta)\), while the information loss \(\varepsilon_{\PE}\) involves the conditional distribution \(P_{Y|X,P}\) and the KL divergence caused by the encoding. From kernel approximation error to information loss requires at least two Pinsker applications and assumptions about the smoothness of \(P_{Y|X,P}\). These assumptions are not stated. The exponential decay function form \(\exp(-d/d_0)\) is guess-level precision.

\end{document}
